\scanpage{40}
% Top note: horizontal separating stroke with a stepped return toward the starred insertion.
\margin{\(^*\) It suffices for this to hold that
\(\delta=d\cap d'\) be an initial segment of \(d'\) and contain
the greatest element of \(d\), which \unsure{says that for the|says that by the}
\(\textcircled{D}\), \(\Sigma_d-\Sigma_\delta\) is
monomorphic and \(\Sigma_{d''}-\Sigma_d\) is fibrant}

(3) Let \(d,d'\) be two flags such that \inserted{simply\(^*\)}
the last element \inserted{\(i\)} of \(d\) is the first element of \(d'\)
and consider \(d''=d\cup d'\) and the diagram
% The extra upper-right arrow has a looped/retraced tail and an unwritten target.
\diagram{\begin{tikzcd}
 \Sigma_{d''}\arrow[r,hook]\arrow[d]&\Sigma_{d'}\arrow[d]\arrow[r,hook]&{}\\
 \Sigma_d\arrow[r,hook]&\Sigma_i&
\end{tikzcd}}
the latter
is \uline{cartesian}. (For \(d\) or \(d'\)
of length \(0\) i.e. reduced : \(\{i\}\), this condition
is trivial)

% Right marginal note: enclosing right/bottom stroke and a left-pointing attachment leader.
\margin{It suffices to assume it for
\(d\) of length \(2\) (or \(d'\) of length \(2\)).
One deduces more from this
\deleted{\unsure{the generality that if}}
\inserted{moreover}
\deleted{\unsure{arbitrary}}
\(d,d'\) are two flags such that
\inserted{\(d''=d\cup d'\)}
is a flag and that
\((\delta=d\cap d')\)
\(\deleted{d\cup d^{\unsure{\prime|\prime\prime}}}\)
\(d''=d\amalg_\delta d'\)
cat. or an ordered \deleted{\unsure{set}},
\(^*\) then
\[
 \textcircled{D}
\]
\diagram{\begin{tikzcd}
 \Sigma_{d''}\arrow[r]\arrow[d]&\Sigma_{d'}\arrow[d]\\
 \Sigma_d\arrow[r]&\Sigma_\delta
\end{tikzcd}}
is cartesian.}

\margin{Moreover, one must assume that when
\(d\) is of length \(2\), then
\([\Sigma_{d'}\longrightarrow\Sigma_\delta]\) is
\uline{transverse} to the cylindrical subspaces \(\Sigma_d\) of
\(\Sigma_\delta\) --- and likewise for the cylindrical structures}

% Long left-margin addition, written diagonally; its attachment arrows are not mathematical map arrows.
\margin{(\unsure{as in the preceding number})
One assumes
further that \deleted{\unsure{besides}} if for \(i\in I\),
\(\operatorname{Dr}(I,i)\) denotes the ordered subset
of \(\operatorname{Dr}(I)\) consisting of the flags
\(d\) whose last element is \(i\),
so that the \(\Sigma_d\) are \(\subset\Sigma_i\),
then the \(\Sigma_d\) form in \(\Sigma_i\) a
\uline{precylindrical stratification} of
type \(\operatorname{Dr}(I,i)^\circ\), which is expressed by
the additional conditions
\[
 \text{a)}\quad \Sigma_{kj}\cap\Sigma_{\ell j}=\phi
 \quad\text{if }k,\ell\le j\text{ are incomparable}.
\]
\[
 \text{b)}\quad
 \Sigma_{i_0-i_n}
 =\Sigma_{i_0i_n}\cap\Sigma_{i_1i_n}-\, -
 \cap\Sigma_{\unsure{i_{n-1}i_n}}
 \qquad\text{in }\Sigma_{i_n}
\]
\[
 \left(\Longleftrightarrow\
 \Sigma_d\cap\Sigma_{d'}=
 \begin{cases}
 \phi&\text{if }d\cup d'\text{ not a flag}\\
 \Sigma_{d\cup d'}&\text{otherwise}
 \end{cases}\right)
\]
c) The intersection in b) is transverse.
One assumes more precisely, \unsure{chosen|chosen} cylindrical
structures \unsure{which make transverse}
\(\Sigma_{ij}\) \(\unsure{(i<j)}\).}

This gives us a \unsure{canonical} procedure for
reconstructing the \(\Sigma_d\), up to \inserted{canonical} isomorphism,
when one knows the \(\Sigma_i,\Sigma_{ij}\)
\inserted{\((i<j)\)} and the morphisms
\diagram{\begin{tikzcd}
 \Sigma_{ij}\arrow[r,hook]\arrow[d]&\Sigma_j\\
 \Sigma_i&
\end{tikzcd}}
\deleted{For}
(NB \unsure{that|which}
the vertical arrow is fibrant, and the horizontal one
a boundary morphism). To have the
transition morphisms between the
\(\Sigma_d\) (constructed by induction on \(\operatorname{long}d\),
the \(0\) and \(1\) being already given) one must give oneself
the morphisms
\[
 \underbrace{\Sigma_{ij}\mathbin{\times_{\Sigma_j}}\Sigma_{jk}}
 _{\mathrel{\mathop=^{\mathrm{def}}}\Sigma_{ij\unsure{k|n}}}
 \longrightarrow\Sigma_{i,k}
\]
satisfying the conditions

\margin{39}
